\documentclass[11pt,a4paper,oneside]{book}

\usepackage{fontspec}
\usepackage[brazilian]{babel}
\usepackage{geometry}
\usepackage{xcolor}
\usepackage{hyperref}
\usepackage{bookmark}
\usepackage{enumitem}
\usepackage{booktabs}
\usepackage{longtable}
\usepackage{array}
\usepackage{tabularx}
\usepackage{fancyhdr}
\usepackage{titlesec}
\usepackage{listings}
\usepackage[most]{tcolorbox}
\usepackage{caption}
\usepackage{amsmath}
\usepackage{amssymb}
\usepackage{microtype}
\usepackage{fvextra}
\usepackage{dirtree}

\geometry{margin=2.1cm}
\setmainfont{DejaVu Serif}
\setsansfont{DejaVu Sans}
\setmonofont{DejaVu Sans Mono}

\definecolor{DevButterGreen}{HTML}{39FF88}
\definecolor{DevButterPink}{HTML}{FF4FD8}
\definecolor{DevButterDark}{HTML}{111318}
\definecolor{DevButterPanel}{HTML}{1A1D24}
\definecolor{DevButterSoft}{HTML}{EDFDF4}
\definecolor{CodeBg}{HTML}{0F1117}
\definecolor{CodeFrame}{HTML}{2F3644}
\definecolor{CodeText}{HTML}{E6E6E6}
\definecolor{CodeKeyword}{HTML}{7DD3FC}
\definecolor{CodeString}{HTML}{FDA4AF}
\definecolor{CodeComment}{HTML}{A7F3D0}
\definecolor{CodeNumber}{HTML}{C4B5FD}

\hypersetup{
  colorlinks=true,
  linkcolor=DevButterPink,
  urlcolor=DevButterPink,
  citecolor=DevButterPink,
  pdftitle={Apple Health Export: Projeto de Parsing, Performance e Engenharia em Python},
  pdfauthor={Marcos Pereira / DevButter}
}

\pagestyle{fancy}
\fancyhf{}
\lhead{Apple Health Export}
\rhead{Python, Postgres e Performance}
\cfoot{\thepage}
\setlength{\headheight}{14pt}

\titleformat{\chapter}[display]
  {\normalfont\bfseries\Huge}
  {\color{DevButterPink}\chaptertitlename\ \thechapter}
  {10pt}
  {\color{DevButterDark}}

\titleformat{\section}
  {\normalfont\bfseries\Large\color{DevButterDark}}
  {\thesection}{0.6em}{}

\titleformat{\subsection}
  {\normalfont\bfseries\large\color{DevButterDark}}
  {\thesubsection}{0.6em}{}

\lstdefinestyle{devbutterpython}{
  language=Python,
  basicstyle=\ttfamily\small\color{CodeText},
  backgroundcolor=\color{CodeBg},
  keywordstyle=\color{CodeKeyword}\bfseries,
  stringstyle=\color{CodeString},
  commentstyle=\color{CodeComment}\itshape,
  numberstyle=\tiny\color{CodeNumber},
  numbers=left,
  stepnumber=1,
  numbersep=8pt,
  showstringspaces=false,
  breaklines=true,
  breakatwhitespace=false,
  frame=single,
  rulecolor=\color{CodeFrame},
  tabsize=4,
  upquote=true,
  columns=fullflexible,
  keepspaces=true
}

\lstdefinestyle{devbuttersql}{
  language=SQL,
  basicstyle=\ttfamily\small\color{CodeText},
  backgroundcolor=\color{CodeBg},
  keywordstyle=\color{CodeKeyword}\bfseries,
  stringstyle=\color{CodeString},
  commentstyle=\color{CodeComment}\itshape,
  numberstyle=\tiny\color{CodeNumber},
  numbers=left,
  stepnumber=1,
  numbersep=8pt,
  showstringspaces=false,
  breaklines=true,
  frame=single,
  rulecolor=\color{CodeFrame},
  tabsize=4,
  upquote=true,
  columns=fullflexible,
  keepspaces=true
}

\lstdefinestyle{devbutterbash}{
  language=bash,
  basicstyle=\ttfamily\small\color{CodeText},
  backgroundcolor=\color{CodeBg},
  keywordstyle=\color{CodeKeyword}\bfseries,
  stringstyle=\color{CodeString},
  commentstyle=\color{CodeComment}\itshape,
  numberstyle=\tiny\color{CodeNumber},
  numbers=left,
  stepnumber=1,
  numbersep=8pt,
  showstringspaces=false,
  breaklines=true,
  frame=single,
  rulecolor=\color{CodeFrame},
  tabsize=4,
  upquote=true,
  columns=fullflexible,
  keepspaces=true
}

\newtcbox{\conceptbox}{
  on line,
  colback=DevButterSoft,
  colframe=DevButterGreen!70!black,
  boxrule=0.4pt,
  arc=2pt,
  left=3pt,
  right=3pt,
  top=2pt,
  bottom=2pt
}

\newtcolorbox{studybox}[1]{
  colback=DevButterSoft,
  colframe=DevButterGreen!65!black,
  title={#1},
  fonttitle=\bfseries,
  arc=3pt,
  boxrule=0.6pt
}

\newtcolorbox{warningbox}[1]{
  colback=DevButterPink!6!white,
  colframe=DevButterPink!75!black,
  title={#1},
  fonttitle=\bfseries,
  arc=3pt,
  boxrule=0.6pt
}

\newcommand{\code}[1]{\texttt{\detokenize{#1}}}

\title{\Huge Apple Health Export\\[0.3em]\Large Projeto de Parsing, Performance e Engenharia em Python}
\author{Marcos Pereira / DevButter}
\date{2026}

\begin{document}
\frontmatter
\maketitle
\tableofcontents

\mainmatter
\chapter{Introducao: o que este projeto quer resolver}

A ideia central do projeto e criar um aplicativo capaz de ler os dados exportados do Apple Health no iPhone e transformar esses arquivos em uma base estruturada para gerar insights pessoais. O ponto importante e que o arquivo exportado pode ser grande: centenas de megabytes de XML e varios arquivos GPX. Portanto, este nao e apenas um exercicio de leitura de arquivo. E um problema real de engenharia de dados local.

\begin{studybox}{Objetivo principal}
Criar uma pipeline que leia o Apple Health Export, processe os dados sem explodir memoria, salve em Postgres e permita consultas para gerar insights como batimento cardiaco medio, passos por dia, sono, treinos, calorias, distancia e tendencias ao longo do tempo.
\end{studybox}

O fluxo do projeto sera:

\begin{verbatim}
Apple Health export bruto
        -> streaming parser
        -> fila / batches
        -> Postgres
        -> queries otimizadas
        -> insights
        -> dashboard / app
\end{verbatim}

A primeira decisao importante e: nao analisar o XML gigante diretamente toda vez. O XML e o dado bruto. Ele deve ser parseado uma vez para uma base estruturada. Depois disso, os insights devem consultar o banco.

\section{Por que este projeto e bom para estudo?}

Porque ele junta varios conceitos que normalmente aparecem separados:

\begin{itemize}[leftmargin=1.2cm]
  \item arquivos grandes e \conceptbox{streaming};
  \item escrita em banco e \conceptbox{I/O-bound};
  \item parsing e conversoes com parte \conceptbox{CPU-bound};
  \item \conceptbox{filas}, \conceptbox{threads}, \conceptbox{multiprocessing} e \conceptbox{async/await};
  \item \conceptbox{dataclass}, \conceptbox{batching}, \conceptbox{logging}, \conceptbox{profiling};
  \item estruturas de dados como \conceptbox{list}, \conceptbox{dict}, \conceptbox{set}, \conceptbox{queue};
  \item conceitos de Big O: \conceptbox{O(1)}, \conceptbox{O(n)}, \conceptbox{O(log n)}, \conceptbox{O(n²)};
  \item Postgres, indices, \code{EXPLAIN ANALYZE} e performance de queries.
\end{itemize}

Ou seja: o app vira um desafio pratico de engenharia em Python.

\section{O que nao fazer}

Evite carregar o XML inteiro em memoria:

\begin{lstlisting}[style=devbutterpython]
content = file.read_text(encoding="utf-8")
soup = BeautifulSoup(content, "xml")
\end{lstlisting}

Essa abordagem pode funcionar para XML pequeno, mas e ruim para arquivos de 280 MB, 628 MB ou maiores. Ela carrega todo o texto e ainda monta uma arvore completa em memoria.

\section{O que fazer}

Usar parsing incremental:

\begin{lstlisting}[style=devbutterpython]
import xml.etree.ElementTree as ET

for event, element in ET.iterparse(xml_file, events=("end",)):
    # processa o elemento
    element.clear()
\end{lstlisting}

Esse modelo le um pedaco, processa e limpa. A ideia e manter memoria quase constante.

\section{Meta de aprendizado}

A regra do projeto sera:

\begin{center}
\Large primeiro funciona $\rightarrow$ depois mede $\rightarrow$ depois otimiza $\rightarrow$ depois compara
\end{center}

Isso impede que voce comece com async, multiprocessing, threads, Postgres, indices e dashboard tudo ao mesmo tempo. Primeiro construiremos uma versao simples e correta. Depois adicionaremos complexidade de forma controlada.

\chapter{Mapa dos conceitos que o projeto vai treinar}

Este capitulo apresenta cada conceito rapidamente e mostra onde ele aparece no projeto. A ideia e voce reconhecer o nome, entender a intuicao e depois aplicar no codigo.

\section{I/O-bound}

I/O significa \textit{Input/Output}. Um trecho e I/O-bound quando o programa passa boa parte do tempo esperando algo externo a CPU.

Exemplos:

\begin{itemize}
  \item ler arquivo do disco;
  \item escrever arquivo;
  \item inserir dados no Postgres;
  \item consultar banco;
  \item fazer request HTTP;
  \item esperar uma resposta de rede.
\end{itemize}

No projeto:

\begin{verbatim}
ler export.xml          -> I/O
inserir no Postgres     -> I/O
ler arquivos .gpx       -> I/O
salvar logs             -> I/O
\end{verbatim}

\section{CPU-bound}

Um trecho e CPU-bound quando o gargalo e computacao. O programa esta gastando tempo calculando, convertendo ou processando dados.

No projeto:

\begin{verbatim}
parsear XML em objetos Python       -> parcialmente CPU-bound
converter strings para datas        -> CPU-bound leve
calcular features e estatisticas    -> CPU-bound
normalizar milhoes de pontos GPS    -> CPU-bound
\end{verbatim}

\section{GIL}

GIL significa \textit{Global Interpreter Lock}. No CPython, ele impede que varias threads executem bytecode Python puro ao mesmo tempo no mesmo processo.

Consequencia pratica:

\begin{itemize}
  \item Threads ajudam bastante quando o gargalo e I/O.
  \item Threads geralmente nao aceleram muito trabalho CPU-bound puro.
  \item Multiprocessing ajuda CPU-bound porque cada processo tem seu proprio interpretador e seu proprio GIL.
  \item Async ajuda I/O concorrente, mas nao deixa calculo CPU-bound magicamente paralelo.
\end{itemize}

\section{Stream}

Stream e processar dados aos poucos. Em vez de carregar tudo, voce consome item por item.

\begin{lstlisting}[style=devbutterpython]
for record in stream_records(xml_file):
    process(record)
\end{lstlisting}

No projeto, \code{ET.iterparse} sera a base do streaming.

\section{Fila}

Fila e uma estrutura FIFO: \textit{first in, first out}. O primeiro item que entra e o primeiro que sai.

No projeto:

\begin{verbatim}
producer/parser -> Queue -> consumer/writer Postgres
\end{verbatim}

A fila desacopla velocidades diferentes. Se o parser for mais rapido que o banco, a fila absorve temporariamente. Se a fila encher, o parser espera. Isso e \textbf{backpressure}.

\section{Batching}

Batching e agrupar varias operacoes em uma so.

Ruim:

\begin{verbatim}
1 registro -> INSERT
1 registro -> INSERT
1 registro -> INSERT
\end{verbatim}

Melhor:

\begin{verbatim}
5000 registros -> INSERT batch
5000 registros -> INSERT batch
5000 registros -> INSERT batch
\end{verbatim}

O algoritmo ainda e O(n), mas o fator constante melhora muito.

\section{Logging}

Logging e registrar eventos importantes do programa. Neste projeto, logs devem mostrar:

\begin{itemize}
  \item quantos registros foram processados;
  \item tamanho do batch;
  \item tamanho da fila;
  \item tempo por insert;
  \item registros por segundo;
  \item erros de parsing ou banco.
\end{itemize}

\section{Profiling}

Profiling e medir onde o programa gasta tempo. Sem profiling, otimizacao vira chute.

Medidas importantes:

\begin{itemize}
  \item tempo total;
  \item tempo por batch;
  \item registros por segundo;
  \item uso de memoria;
  \item query time no banco.
\end{itemize}

\section{Big O aplicado}

\begin{longtable}{p{3cm}p{10cm}}
\toprule
\textbf{Complexidade} & \textbf{Aplicacao no projeto} \\
\midrule
O(1) & acesso medio em \code{dict}, \code{set}, append em lista \\
O(n) & percorrer todos os elementos do XML uma vez \\
O(log n) & busca usando indice B-tree no Postgres \\
O(n log n) & ordenacao de registros por data \\
O(n²) & comparacoes ingênuas entre todos os pares de duas listas; algo a evitar \\
\bottomrule
\end{longtable}

\section{Resumo do mapa}

\begin{longtable}{p{4cm}p{9cm}}
\toprule
\textbf{Conceito} & \textbf{Onde aparece} \\
\midrule
list / array & batches de registros \\
dict / hash map & metadata, contadores, mapeamento de tipos \\
set / hash set & filtro rapido de tipos aceitos \\
queue & producer/consumer \\
tree & XML e uma arvore \\
stack & percursos DFS opcionais \\
recursion & exploracao de XML pequeno; evitar para XML gigante \\
sorting & ordenar por data \\
binary search & achar registros em intervalos ordenados \\
two pointers & cruzar intervalos de sono com batimentos \\
sliding window & medias moveis de 7/14 dias \\
caching & summary tables e materialized views \\
indexing & indices por tipo e data no Postgres \\
\bottomrule
\end{longtable}

\chapter{Estrutura do projeto}

Voce ja tem uma estrutura inicial. A ideia e manter essa base e agregar os modulos necessarios aos poucos.

\section{Estrutura inicial fornecida}

\begin{verbatim}
apple_health_export
├─ .editorconfig
├─ README.md
├─ notebooks
│  ├─ README.md
│  └─ playground.ipynb
├─ poetry.lock
├─ pyproject.toml
└─ src
   ├─ __init__.py
   ├─ app.py
   ├─ config
   │  ├─ __init__.py
   │  ├─ logger.py
   │  └─ settings.py
   └─ parser
      ├─ __init__.py
      ├─ base_parser.py
      └─ xml_parser.py
\end{verbatim}

Essa estrutura ja tem tres decisoes boas:

\begin{itemize}
  \item \code{src/app.py}: ponto de entrada;
  \item \code{src/config}: configuracao e logging;
  \item \code{src/parser}: responsabilidade de parsing separada.
\end{itemize}

\section{Estrutura agregada proposta}

Agora vamos expandir sem destruir sua organizacao:

\begin{verbatim}
apple_health_export
├─ .editorconfig
├─ README.md
├─ docker-compose.yml
├─ .env.example
├─ notebooks
│  ├─ README.md
│  └─ playground.ipynb
├─ poetry.lock
├─ pyproject.toml
├─ tests
│  ├─ __init__.py
│  ├─ test_parser_utils.py
│  ├─ test_xml_parser.py
│  └─ fixtures
│     └─ small_export.xml
└─ src
   ├─ __init__.py
   ├─ app.py
   ├─ config
   │  ├─ __init__.py
   │  ├─ logger.py
   │  └─ settings.py
   ├─ db
   │  ├─ __init__.py
   │  ├─ connection.py
   │  ├─ schema.py
   │  └─ repository.py
   ├─ models
   │  ├─ __init__.py
   │  └─ health_record.py
   ├─ parser
   │  ├─ __init__.py
   │  ├─ base_parser.py
   │  ├─ xml_parser.py
   │  ├─ gpx_parser.py
   │  └─ profiler.py
   ├─ pipelines
   │  ├─ __init__.py
   │  ├─ sync_pipeline.py
   │  ├─ threaded_pipeline.py
   │  └─ multiprocessing_pipeline.py
   └─ insights
      ├─ __init__.py
      ├─ daily_summary.py
      └─ queries.py
\end{verbatim}

\section{Por que cada pasta existe?}

\begin{longtable}{p{4cm}p{9cm}}
\toprule
\textbf{Pasta/arquivo} & \textbf{Responsabilidade} \\
\midrule
\code{config/settings.py} & centralizar variaveis como host do banco, porta, usuario, senha, batch size \\
\code{config/logger.py} & configurar logs padronizados \\
\code{models/} & dataclasses e objetos de dominio \\
\code{parser/} & transformar XML/GPX em objetos Python \\
\code{db/} & conexao, schema e escrita no Postgres \\
\code{pipelines/} & orquestrar parser, fila, batches e writer \\
\code{insights/} & queries e agregacoes para gerar metricas \\
\code{tests/} & testes pequenos sem usar o XML gigante \\
\bottomrule
\end{longtable}

\section{Regra de separacao}

Cada modulo deve ter uma responsabilidade clara:

\begin{itemize}
  \item Parser nao deve saber detalhes do Docker.
  \item Repository nao deve parsear XML.
  \item Logger nao deve conter regra de negocio.
  \item Pipeline orquestra componentes, mas nao deve concentrar toda logica.
\end{itemize}

Essa separacao facilita testes e torna a evolucao mais segura.

\chapter{I/O e streaming: lendo arquivos grandes sem sofrer}

\section{O que e I/O?}

I/O significa \textit{Input/Output}. Em termos simples, e tudo que envolve entrada ou saida de dados fora da CPU. Ler um arquivo, escrever em disco, consultar banco e esperar uma resposta de rede sao exemplos de I/O.

No seu projeto, as operacoes de I/O mais importantes sao:

\begin{itemize}
  \item ler \code{export.xml};
  \item ler arquivos \code{.gpx};
  \item inserir dados no Postgres;
  \item escrever logs;
  \item consultar tabelas de insights.
\end{itemize}

\section{Por que o XML gigante muda a estrategia?}

Se o arquivo tem 280 MB ou 628 MB, carregar tudo em memoria pode ser ruim. Alem do texto bruto, bibliotecas como BeautifulSoup criam objetos Python para representar a arvore XML inteira. Isso pode multiplicar o uso de memoria.

\begin{warningbox}{Evite no arquivo gigante}
\begin{lstlisting}[style=devbutterpython]
content = file.read_text(encoding="utf-8")
soup = BeautifulSoup(content, "xml")
\end{lstlisting}
\end{warningbox}

BeautifulSoup ainda pode ser util para arquivos pequenos, testes e inspecoes pontuais. Mas nao deve ser o parser principal do export gigante.

\section{Streaming com iterparse}

A alternativa e usar \code{xml.etree.ElementTree.iterparse}. Ele percorre o XML elemento por elemento.

\begin{lstlisting}[style=devbutterpython]
from pathlib import Path
from collections.abc import Iterator
import xml.etree.ElementTree as ET


def clean_tag(tag: str) -> str:
    if "}" in tag:
        return tag.split("}", 1)[1]
    return tag


def stream_record_elements(xml_file: Path) -> Iterator[ET.Element]:
    context = ET.iterparse(xml_file, events=("end",))

    for _, element in context:
        tag = clean_tag(element.tag)

        if tag == "Record":
            yield element

        element.clear()
\end{lstlisting}

\section{Explicando cada nomezinho}

\begin{itemize}
  \item \code{Path}: objeto da biblioteca \code{pathlib} para lidar com caminhos de arquivo de forma mais segura que string.
  \item \code{Iterator}: tipo que indica que a funcao retorna valores aos poucos.
  \item \code{ET}: apelido comum para \code{xml.etree.ElementTree}.
  \item \code{iterparse}: parser incremental de XML.
  \item \code{events=("end",)}: processa o elemento quando ele fecha, ou seja, quando ja temos seus atributos e filhos.
  \item \code{yield}: retorna um item sem encerrar a funcao; cria um generator.
  \item \code{element.clear()}: limpa o elemento da memoria depois de processar.
\end{itemize}

\section{Complexidade}

Se o XML tem \(n\) elementos, o parser passa por cada elemento uma vez:

\[
T(n) = O(n)
\]

A memoria tende a ser quase constante:

\[
M(n) \approx O(1)
\]

Isso nao significa literalmente zero crescimento, mas significa que voce nao precisa manter o XML inteiro em memoria.

\section{Primeiro exercicio}

Crie um XML pequeno em \code{tests/fixtures/small_export.xml} e teste se sua stream encontra apenas tags \code{Record}.

\begin{lstlisting}[style=devbutterpython]
def test_stream_record_elements_returns_records_only(tmp_path):
    xml_file = tmp_path / "small_export.xml"
    xml_file.write_text("""
    <HealthData>
        <Record type="HKQuantityTypeIdentifierHeartRate" value="90" />
        <Workout workoutActivityType="HKWorkoutActivityTypeRunning" />
    </HealthData>
    """)

    records = list(stream_record_elements(xml_file))

    assert len(records) == 1
    assert records[0].attrib["type"] == "HKQuantityTypeIdentifierHeartRate"
\end{lstlisting}

\chapter{Dataclass e modelagem de dados}

\section{Por que modelar antes de inserir no banco?}

Ao parsear XML, voce recebe atributos como strings soltas. Se voce passar dicionarios por todo o projeto, o codigo fica fragil:

\begin{lstlisting}[style=devbutterpython]
record["startDate"]
record["value"]
record["sourceName"]
\end{lstlisting}

O problema e que erros de chave aparecem tarde. Alem disso, nao fica claro quais campos existem.

Com \code{dataclass}, voce cria uma estrutura explicita:

\begin{lstlisting}[style=devbutterpython]
record.start_date
record.value_numeric
record.source_name
\end{lstlisting}

\section{Modelo HealthRecord}

Arquivo sugerido: \code{src/models/health_record.py}

\begin{lstlisting}[style=devbutterpython]
from dataclasses import dataclass


@dataclass(slots=True)
class HealthRecord:
    apple_record_id: int
    type: str | None
    source_name: str | None
    source_version: str | None
    device: str | None
    unit: str | None
    creation_date: str | None
    start_date: str | None
    end_date: str | None
    value: str | None
    value_numeric: float | None
    metadata: dict[str, str]
\end{lstlisting}

\section{Explicando cada nomezinho}

\begin{itemize}
  \item \code{@dataclass}: decorator que gera automaticamente metodos como \code{__init__} e \code{__repr__}.
  \item \code{slots=True}: reduz overhead de memoria removendo o \code{__dict__} dinamico de cada instancia.
  \item \code{apple_record_id}: id sequencial criado pelo seu parser, diferente do id do banco.
  \item \code{type}: tipo do Record do Apple Health, por exemplo \code{HKQuantityTypeIdentifierHeartRate}.
  \item \code{source_name}: origem do dado, por exemplo Apple Watch ou iPhone.
  \item \code{value}: valor original em texto.
  \item \code{value_numeric}: tentativa de conversao para numero.
  \item \code{metadata}: informacoes extras de \code{MetadataEntry}.
\end{itemize}

\section{Por que manter value e value\_numeric?}

Nem todo valor do Apple Health e numerico. Alguns sao categorias.

Numericos:

\begin{verbatim}
92
1200
75.4
\end{verbatim}

Categoricos:

\begin{verbatim}
HKCategoryValueSleepAnalysisAsleepCore
HKCategoryValueSleepAnalysisInBed
\end{verbatim}

Por isso salvamos os dois:

\begin{itemize}
  \item \code{value}: preserva o valor original;
  \item \code{value_numeric}: permite agregacoes numericas quando possivel.
\end{itemize}

\section{Funcao safe\_float}

\begin{lstlisting}[style=devbutterpython]
def safe_float(value: str | None) -> float | None:
    if value is None:
        return None

    try:
        return float(value)
    except ValueError:
        return None
\end{lstlisting}

Ela tenta converter texto para numero. Se nao conseguir, retorna \code{None}. Isso evita quebrar o parser quando encontrar valores categoricos.

\section{Exercicio}

Implemente testes:

\begin{lstlisting}[style=devbutterpython]
def test_safe_float_with_number():
    assert safe_float("92.5") == 92.5


def test_safe_float_with_none():
    assert safe_float(None) is None


def test_safe_float_with_text():
    assert safe_float("HKCategoryValueSleepAnalysisAsleepCore") is None
\end{lstlisting}

\chapter{Postgres desde o inicio}

\section{Por que Postgres em vez de CSV?}

CSV e simples para debugar, mas o objetivo deste projeto tambem e treinar banco, indices, queries e performance. Entao faz sentido usar Postgres desde o inicio.

Com Postgres voce treina:

\begin{itemize}
  \item schema design;
  \item tipos de dados;
  \item \code{JSONB};
  \item batch insert;
  \item transacoes;
  \item indices;
  \item \code{EXPLAIN ANALYZE};
  \item queries de insight.
\end{itemize}

\section{Docker Compose para Postgres}

Arquivo: \code{docker-compose.yml}

\begin{lstlisting}[style=devbutterbash]
services:
  postgres:
    image: postgres:16
    container_name: apple_health_postgres
    environment:
      POSTGRES_USER: apple
      POSTGRES_PASSWORD: apple
      POSTGRES_DB: apple_health
    ports:
      - "5433:5432"
    volumes:
      - postgres_data:/var/lib/postgresql/data

volumes:
  postgres_data:
\end{lstlisting}

Subir banco:

\begin{lstlisting}[style=devbutterbash]
docker compose up -d
\end{lstlisting}

\section{Settings}

Arquivo: \code{src/config/settings.py}

\begin{lstlisting}[style=devbutterpython]
from dataclasses import dataclass
import os


@dataclass(frozen=True, slots=True)
class Settings:
    postgres_host: str = os.getenv("POSTGRES_HOST", "localhost")
    postgres_port: int = int(os.getenv("POSTGRES_PORT", "5433"))
    postgres_user: str = os.getenv("POSTGRES_USER", "apple")
    postgres_password: str = os.getenv("POSTGRES_PASSWORD", "apple")
    postgres_db: str = os.getenv("POSTGRES_DB", "apple_health")
    batch_size: int = int(os.getenv("BATCH_SIZE", "5000"))

    @property
    def database_url(self) -> str:
        return (
            f"postgresql://{self.postgres_user}:"
            f"{self.postgres_password}@"
            f"{self.postgres_host}:{self.postgres_port}/"
            f"{self.postgres_db}"
        )
\end{lstlisting}

\section{Schema inicial}

Tabela principal: \code{health_records}.

\begin{lstlisting}[style=devbuttersql]
CREATE TABLE IF NOT EXISTS health_records (
    id BIGSERIAL PRIMARY KEY,
    apple_record_id BIGINT NOT NULL,
    type TEXT,
    source_name TEXT,
    source_version TEXT,
    device TEXT,
    unit TEXT,
    creation_date TIMESTAMPTZ,
    start_date TIMESTAMPTZ,
    end_date TIMESTAMPTZ,
    value TEXT,
    value_numeric DOUBLE PRECISION,
    metadata JSONB
);
\end{lstlisting}

\section{Tabela parse\_runs}

Essa tabela guarda historico de execucoes do parser.

\begin{lstlisting}[style=devbuttersql]
CREATE TABLE IF NOT EXISTS parse_runs (
    id BIGSERIAL PRIMARY KEY,
    started_at TIMESTAMPTZ NOT NULL DEFAULT NOW(),
    finished_at TIMESTAMPTZ,
    file_path TEXT NOT NULL,
    total_records BIGINT DEFAULT 0,
    status TEXT NOT NULL DEFAULT 'running',
    error TEXT
);
\end{lstlisting}

\section{Indices}

Consultas de insights vao filtrar muito por tipo e data. Entao os primeiros indices sao:

\begin{lstlisting}[style=devbuttersql]
CREATE INDEX IF NOT EXISTS idx_health_records_type
ON health_records (type);

CREATE INDEX IF NOT EXISTS idx_health_records_start_date
ON health_records (start_date);

CREATE INDEX IF NOT EXISTS idx_health_records_type_start_date
ON health_records (type, start_date);
\end{lstlisting}

\section{Por que indice importa?}

Sem indice, o banco pode precisar varrer muitas linhas:

\[
O(n)
\]

Com indice B-tree, buscas por tipo/data podem se aproximar de:

\[
O(log n)
\]

Na pratica, confirme com:

\begin{lstlisting}[style=devbuttersql]
EXPLAIN ANALYZE
SELECT
    DATE(start_date) AS day,
    AVG(value_numeric) AS avg_heart_rate
FROM health_records
WHERE type = 'HKQuantityTypeIdentifierHeartRate'
GROUP BY DATE(start_date)
ORDER BY day;
\end{lstlisting}

\chapter{Parser streaming + batching}

\section{Fluxo desta etapa}

Agora juntamos streaming, dataclass e Postgres.

\begin{verbatim}
XML element
    -> parse attributes
    -> HealthRecord dataclass
    -> batch list
    -> Postgres
\end{verbatim}

\section{Parser de elemento Record}

Arquivo: \code{src/parser/xml_parser.py}

\begin{lstlisting}[style=devbutterpython]
import xml.etree.ElementTree as ET
from src.models.health_record import HealthRecord


def clean_tag(tag: str) -> str:
    if "}" in tag:
        return tag.split("}", 1)[1]
    return tag


def safe_float(value: str | None) -> float | None:
    if value is None:
        return None

    try:
        return float(value)
    except ValueError:
        return None


def extract_metadata(element: ET.Element) -> dict[str, str]:
    metadata: dict[str, str] = {}

    for child in element:
        if clean_tag(child.tag) != "MetadataEntry":
            continue

        key = child.attrib.get("key")
        value = child.attrib.get("value")

        if key is not None and value is not None:
            metadata[key] = value

    return metadata


def parse_record_element(
    element: ET.Element,
    record_id: int,
) -> HealthRecord:
    raw_value = element.attrib.get("value")

    return HealthRecord(
        apple_record_id=record_id,
        type=element.attrib.get("type"),
        source_name=element.attrib.get("sourceName"),
        source_version=element.attrib.get("sourceVersion"),
        device=element.attrib.get("device"),
        unit=element.attrib.get("unit"),
        creation_date=element.attrib.get("creationDate"),
        start_date=element.attrib.get("startDate"),
        end_date=element.attrib.get("endDate"),
        value=raw_value,
        value_numeric=safe_float(raw_value),
        metadata=extract_metadata(element),
    )
\end{lstlisting}

\section{Repository com batch insert}

Instale o driver:

\begin{lstlisting}[style=devbutterbash]
poetry add "psycopg[binary]"
\end{lstlisting}

Arquivo: \code{src/db/repository.py}

\begin{lstlisting}[style=devbutterpython]
import json
from collections.abc import Sequence

import psycopg

from src.models.health_record import HealthRecord


class HealthRecordRepository:
    def __init__(self, connection: psycopg.Connection) -> None:
        self.connection = connection

    def insert_batch(self, records: Sequence[HealthRecord]) -> None:
        if not records:
            return

        rows = [
            (
                record.apple_record_id,
                record.type,
                record.source_name,
                record.source_version,
                record.device,
                record.unit,
                record.creation_date,
                record.start_date,
                record.end_date,
                record.value,
                record.value_numeric,
                json.dumps(record.metadata),
            )
            for record in records
        ]

        with self.connection.cursor() as cursor:
            cursor.executemany(
                """
                INSERT INTO health_records (
                    apple_record_id,
                    type,
                    source_name,
                    source_version,
                    device,
                    unit,
                    creation_date,
                    start_date,
                    end_date,
                    value,
                    value_numeric,
                    metadata
                )
                VALUES (
                    %s, %s, %s, %s, %s, %s,
                    %s, %s, %s, %s, %s, %s
                )
                """,
                rows,
            )

        self.connection.commit()
\end{lstlisting}

\section{Pipeline sincrona}

Arquivo: \code{src/pipelines/sync_pipeline.py}

\begin{lstlisting}[style=devbutterpython]
from pathlib import Path
from time import perf_counter
import logging
import xml.etree.ElementTree as ET

from src.db.repository import HealthRecordRepository
from src.parser.xml_parser import clean_tag, parse_record_element

logger = logging.getLogger(__name__)


def run_sync_pipeline(
    xml_file: Path,
    repository: HealthRecordRepository,
    batch_size: int,
) -> None:
    batch = []
    record_id = 0
    started_at = perf_counter()

    context = ET.iterparse(xml_file, events=("end",))

    for _, element in context:
        tag = clean_tag(element.tag)

        if tag == "Record":
            record_id += 1
            batch.append(parse_record_element(element, record_id))

            if len(batch) >= batch_size:
                insert_started_at = perf_counter()
                repository.insert_batch(batch)
                insert_elapsed = perf_counter() - insert_started_at

                logger.info(
                    "Inserted %s records in %.2fs | total=%s",
                    len(batch),
                    insert_elapsed,
                    record_id,
                )

                batch.clear()

        element.clear()

    if batch:
        repository.insert_batch(batch)

    elapsed = perf_counter() - started_at
    logger.info(
        "Finished parsing %s records in %.2fs | %.2f records/s",
        record_id,
        elapsed,
        record_id / elapsed if elapsed > 0 else 0,
    )
\end{lstlisting}

\section{Conceitos aplicados}

\begin{itemize}
  \item \code{batch}: uma \code{list}; append e O(1) medio.
  \item \code{batch_size}: tradeoff entre memoria e velocidade.
  \item \code{iterparse}: lazy loading / streaming.
  \item \code{insert_batch}: reduz overhead de I/O com o banco.
  \item \code{perf_counter}: mede tempo com precisao suficiente para benchmark.
\end{itemize}

\section{Experimento}

Rode a pipeline com:

\begin{verbatim}
batch_size = 500
batch_size = 5000
batch_size = 20000
batch_size = 50000
\end{verbatim}

Anote:

\begin{itemize}
  \item tempo total;
  \item records/s;
  \item memoria percebida;
  \item tempo medio por batch.
\end{itemize}

\chapter{Filas e threads: producer/consumer}

\section{Por que usar fila?}

Parser e banco podem ter velocidades diferentes. O parser pode produzir registros mais rapido do que o Postgres consegue inserir, ou o banco pode ficar esperando I/O enquanto o parser poderia continuar trabalhando.

A arquitetura producer/consumer separa essas responsabilidades:

\begin{verbatim}
producer thread
    le XML e cria HealthRecord
        -> Queue
consumer thread
    pega HealthRecord e insere em batch no Postgres
\end{verbatim}

\section{Queue}

\code{queue.Queue} e uma fila thread-safe da biblioteca padrao.

\begin{lstlisting}[style=devbutterpython]
from queue import Queue

record_queue: Queue[HealthRecord | None] = Queue(maxsize=10_000)
\end{lstlisting}

\section{Explicando cada nomezinho}

\begin{itemize}
  \item \code{Queue}: fila segura para multiplas threads.
  \item \code{maxsize}: limite da fila. Evita crescimento infinito de memoria.
  \item \code{HealthRecord | None}: a fila aceita registros ou \code{None}.
  \item \code{None}: usado como sentinela para sinalizar fim do trabalho.
  \item \code{backpressure}: quando a fila enche, o producer bloqueia. Isso protege a memoria.
\end{itemize}

\section{Producer}

\begin{lstlisting}[style=devbutterpython]
from pathlib import Path
from queue import Queue
import xml.etree.ElementTree as ET

from src.models.health_record import HealthRecord
from src.parser.xml_parser import clean_tag, parse_record_element


def producer(
    xml_file: Path,
    record_queue: Queue[HealthRecord | None],
) -> None:
    record_id = 0
    context = ET.iterparse(xml_file, events=("end",))

    for _, element in context:
        if clean_tag(element.tag) == "Record":
            record_id += 1
            record = parse_record_element(element, record_id)
            record_queue.put(record)

        element.clear()

    record_queue.put(None)
\end{lstlisting}

\section{Consumer}

\begin{lstlisting}[style=devbutterpython]
from queue import Queue

from src.db.repository import HealthRecordRepository
from src.models.health_record import HealthRecord


def consumer(
    record_queue: Queue[HealthRecord | None],
    repository: HealthRecordRepository,
    batch_size: int,
) -> None:
    batch: list[HealthRecord] = []

    while True:
        record = record_queue.get()

        if record is None:
            break

        batch.append(record)

        if len(batch) >= batch_size:
            repository.insert_batch(batch)
            batch.clear()

    if batch:
        repository.insert_batch(batch)
\end{lstlisting}

\section{Rodando com Thread}

\begin{lstlisting}[style=devbutterpython]
from pathlib import Path
from queue import Queue
from threading import Thread

from src.db.repository import HealthRecordRepository
from src.models.health_record import HealthRecord


def run_threaded_pipeline(
    xml_file: Path,
    repository: HealthRecordRepository,
    batch_size: int,
) -> None:
    record_queue: Queue[HealthRecord | None] = Queue(maxsize=10_000)

    producer_thread = Thread(
        target=producer,
        args=(xml_file, record_queue),
        name="xml-producer",
    )

    consumer_thread = Thread(
        target=consumer,
        args=(record_queue, repository, batch_size),
        name="postgres-consumer",
    )

    producer_thread.start()
    consumer_thread.start()

    producer_thread.join()
    consumer_thread.join()
\end{lstlisting}

\section{Quando thread ajuda?}

Threads ajudam quando existe espera de I/O. No projeto, escrever no Postgres e ler arquivos sao I/O-bound. Enquanto uma thread espera o banco responder, outra pode continuar lendo ou preparando dados.

Mas se o gargalo principal for parse CPU-bound, threads podem nao acelerar muito por causa do GIL.

\section{Experimento}

Compare:

\begin{itemize}
  \item pipeline sincrona;
  \item pipeline com producer/consumer e threads;
  \item batch sizes diferentes;
  \item queue maxsize diferente.
\end{itemize}

Pergunta que voce deve responder: \textbf{thread melhorou por causa de I/O ou nao mudou por causa de CPU/GIL?}

\chapter{GIL, multiprocessing e async/await}

\section{GIL na pratica}

O GIL impede que multiplas threads executem bytecode Python puro simultaneamente no mesmo processo. Isso nao significa que threads sao inuteis. Significa que voce precisa entender o tipo de gargalo.

\begin{longtable}{p{4cm}p{8cm}}
\toprule
\textbf{Cenario} & \textbf{Ferramenta provavel} \\
\midrule
I/O-bound & threads ou async \\
CPU-bound & multiprocessing \\
Muitos arquivos pequenos & ThreadPoolExecutor pode ajudar \\
Transformacao numerica pesada & ProcessPoolExecutor pode ajudar \\
API web com muitas esperas & async/await pode ajudar \\
\bottomrule
\end{longtable}

\section{Multiprocessing}

Multiprocessing cria processos separados. Cada processo tem seu proprio interpretador Python e seu proprio GIL. Por isso ele pode acelerar tarefas CPU-bound.

Exemplo didatico:

\begin{lstlisting}[style=devbutterpython]
from concurrent.futures import ProcessPoolExecutor


def cpu_heavy_transform(value: float | None) -> float | None:
    if value is None:
        return None

    result = value
    for _ in range(100_000):
        result = (result * 1.000001) / 1.000001

    return result


def process_values(values: list[float | None]) -> list[float | None]:
    with ProcessPoolExecutor() as executor:
        return list(executor.map(cpu_heavy_transform, values))
\end{lstlisting}

\section{Custo do multiprocessing}

Multiprocessing nao e gratis:

\begin{itemize}
  \item cada processo consome mais memoria;
  \item dados precisam ser serializados entre processos;
  \item criar processos tem overhead;
  \item para tarefas pequenas, pode ficar mais lento.
\end{itemize}

Use multiprocessing quando o trabalho por item for pesado o suficiente.

\section{ThreadPoolExecutor para GPX}

Arquivos GPX sao bons candidatos para concorrencia porque geralmente existem varios arquivos menores.

\begin{lstlisting}[style=devbutterpython]
from concurrent.futures import ThreadPoolExecutor
from pathlib import Path


def parse_gpx_file(file: Path) -> list[dict[str, str | None]]:
    # parseia um arquivo GPX e retorna pontos
    return []


def parse_many_gpx(files: list[Path]) -> list[dict[str, str | None]]:
    rows: list[dict[str, str | None]] = []

    with ThreadPoolExecutor(max_workers=8) as executor:
        for file_rows in executor.map(parse_gpx_file, files):
            rows.extend(file_rows)

    return rows
\end{lstlisting}

\section{Async/await}

\code{async/await} e uma forma de concorrencia cooperativa. Em vez de criar varias threads, um event loop alterna entre tarefas quando elas fazem \code{await} em operacoes I/O.

Exemplo minimo:

\begin{lstlisting}[style=devbutterpython]
import asyncio


async def producer(queue: asyncio.Queue[int]) -> None:
    for item in range(10):
        await queue.put(item)

    await queue.put(-1)


async def consumer(queue: asyncio.Queue[int]) -> None:
    while True:
        item = await queue.get()

        if item == -1:
            break

        print(item)


async def main() -> None:
    queue: asyncio.Queue[int] = asyncio.Queue()

    await asyncio.gather(
        producer(queue),
        consumer(queue),
    )


asyncio.run(main())
\end{lstlisting}

\section{Onde async entra no projeto?}

Nao comece com async no parser principal. Primeiro faca a versao sincrona eficiente. Async pode entrar depois em:

\begin{itemize}
  \item API com FastAPI;
  \item consulta de status de parse;
  \item jobs em background;
  \item writer async com \code{asyncpg};
  \item dashboard consultando dados sem bloquear.
\end{itemize}

\section{Ordem recomendada}

\begin{enumerate}
  \item parser sincrono;
  \item logs e profiling;
  \item threads com queue;
  \item ThreadPool para GPX;
  \item ProcessPool para CPU-bound;
  \item async para API ou I/O concorrente futuro.
\end{enumerate}

\chapter{Logging e profiling}

\section{Por que logging e essencial?}

Sem logs, voce nao sabe se o parser esta travado, lento, quebrando em algum Record especifico ou inserindo pouco por segundo.

Logs devem responder:

\begin{itemize}
  \item quantos registros ja foram processados?
  \item quantos registros por segundo?
  \item qual o tamanho do batch?
  \item quanto tempo cada insert demorou?
  \item qual thread esta executando?
  \item ocorreu erro em qual arquivo?
\end{itemize}

\section{Configuracao de logger}

Arquivo: \code{src/config/logger.py}

\begin{lstlisting}[style=devbutterpython]
import logging
import sys


def configure_logging(level: int = logging.INFO) -> None:
    logging.basicConfig(
        level=level,
        format=(
            "%(asctime)s | %(levelname)s | "
            "%(threadName)s | %(name)s | %(message)s"
        ),
        handlers=[logging.StreamHandler(sys.stdout)],
    )
\end{lstlisting}

\section{Uso}

\begin{lstlisting}[style=devbutterpython]
import logging

logger = logging.getLogger(__name__)

logger.info("Starting Apple Health parser")
logger.warning("Skipping unsupported tag: %s", tag)
logger.exception("Failed to insert batch")
\end{lstlisting}

\section{Profiling simples com perf\_counter}

\begin{lstlisting}[style=devbutterpython]
from time import perf_counter

started_at = perf_counter()

# operacao que voce quer medir

elapsed = perf_counter() - started_at
print(f"Elapsed: {elapsed:.2f}s")
\end{lstlisting}

\section{Throughput}

Throughput e taxa de processamento. Neste projeto, uma metrica boa e registros por segundo.

\begin{lstlisting}[style=devbutterpython]
records_per_second = total_records / elapsed
\end{lstlisting}

Exemplo de log desejado:

\begin{verbatim}
2026-06-23 08:00:00 | INFO | MainThread | parser | Inserted 5000 records in 0.41s | 12195 records/s
\end{verbatim}

\section{Benchmark de batch size}

Crie uma tabela manual:

\begin{longtable}{rrrr}
\toprule
\textbf{Batch size} & \textbf{Tempo total} & \textbf{Records/s} & \textbf{Observacao} \\
\midrule
500 & & & muitos inserts \\
5000 & & & equilibrio inicial \\
20000 & & & menos commits \\
50000 & & & maior memoria \\
\bottomrule
\end{longtable}

\section{Profiling mais avancado}

Depois voce pode usar:

\begin{lstlisting}[style=devbutterbash]
python -m cProfile -o profile.out src/app.py parse
\end{lstlisting}

E analisar com:

\begin{lstlisting}[style=devbutterbash]
python -m pstats profile.out
\end{lstlisting}

Mas comece com \code{perf_counter}, porque ele te obriga a formular exatamente o que esta medindo.

\chapter{Algoritmos e estruturas de dados aplicados}

\section{list / array}

Em Python, \code{list} e um array dinamico. No projeto, ela aparece como batch:

\begin{lstlisting}[style=devbutterpython]
batch: list[HealthRecord] = []
batch.append(record)
\end{lstlisting}

\code{append} e O(1) medio. Acesso por indice tambem e O(1).

\section{dict / hash map}

\code{dict} e uma tabela hash. Usamos para metadata e contadores.

\begin{lstlisting}[style=devbutterpython]
metadata: dict[str, str] = {}
metadata[key] = value
\end{lstlisting}

Lookup medio:

\[
O(1)
\]

\section{set / hash set}

\code{set} e usado para membership rapido.

\begin{lstlisting}[style=devbutterpython]
allowed_types = {
    "HKQuantityTypeIdentifierHeartRate",
    "HKQuantityTypeIdentifierStepCount",
}

if record.type in allowed_types:
    process(record)
\end{lstlisting}

Checar se um item esta no set e O(1) medio.

\section{queue}

Fila FIFO usada em producer/consumer.

\begin{lstlisting}[style=devbutterpython]
record_queue.put(record)
record = record_queue.get()
\end{lstlisting}

\section{tree}

XML e uma arvore:

\begin{verbatim}
Workout
 ├── MetadataEntry
 ├── WorkoutStatistics
 └── WorkoutEvent
\end{verbatim}

Cada tag e um no; tags internas sao filhos.

\section{recursion}

Recursao pode percorrer arvores pequenas:

\begin{lstlisting}[style=devbutterpython]
def walk(element):
    print(element.tag)

    for child in element:
        walk(child)
\end{lstlisting}

Mas evite usar recursao carregando a arvore inteira do XML gigante. Para o export real, prefira streaming.

\section{sorting}

Ordenar por data e comum em series temporais. Ordenacao geralmente custa:

\[
O(n \log n)
\]

No banco:

\begin{lstlisting}[style=devbuttersql]
SELECT *
FROM health_records
ORDER BY start_date;
\end{lstlisting}

\section{binary search}

Busca binaria encontra um item em lista ordenada em O(log n). Em Python:

\begin{lstlisting}[style=devbutterpython]
from bisect import bisect_left

positions = [1, 3, 5, 7, 9]
index = bisect_left(positions, 6)
\end{lstlisting}

No projeto, a ideia aparece ao buscar registros em series temporais ordenadas. No banco, indices B-tree usam uma logica parecida de busca eficiente.

\section{two pointers}

Dois ponteiros servem para cruzar duas listas ordenadas sem fazer O(n²).

Exemplo: achar batimentos dentro de intervalos de sono.

\begin{verbatim}
sono:       [inicio, fim]
heart rate: timestamp ordenado
\end{verbatim}

Em vez de comparar todo sono com todo batimento, voce avanca ponteiros conforme as datas.

\section{sliding window}

Sliding window e uma janela movel. Exemplo: media movel de 7 dias.

Em SQL:

\begin{lstlisting}[style=devbuttersql]
SELECT
    day,
    steps,
    AVG(steps) OVER (
        ORDER BY day
        ROWS BETWEEN 6 PRECEDING AND CURRENT ROW
    ) AS steps_7d_avg
FROM daily_summary
ORDER BY day;
\end{lstlisting}

Isso e muito util para tendencias de passos, sono e batimento.

\section{O(n²): o que evitar}

Evite loops aninhados sem necessidade:

\begin{lstlisting}[style=devbutterpython]
for sleep_interval in sleep_intervals:
    for heart_rate in heart_rates:
        compare(sleep_interval, heart_rate)
\end{lstlisting}

Se cada lista tem muitos itens, isso vira O(n²). Busque alternativas com ordenacao, indices, SQL, two pointers ou janelas.

\chapter{Testes e exercicios}

\section{Por que testar com XML pequeno?}

Voce nao deve testar com o arquivo gigante. Testes devem ser pequenos, rapidos e previsiveis.

Crie:

\begin{verbatim}
tests/fixtures/small_export.xml
\end{verbatim}

Conteudo exemplo:

\begin{lstlisting}[style=devbutterbash]
<HealthData>
    <ExportDate value="2026-06-23 08:00:00 -0300" />
    <Record
        type="HKQuantityTypeIdentifierHeartRate"
        sourceName="Apple Watch"
        unit="count/min"
        startDate="2026-06-23 07:00:00 -0300"
        endDate="2026-06-23 07:00:00 -0300"
        value="90"
    />
    <Record
        type="HKCategoryTypeIdentifierSleepAnalysis"
        sourceName="Apple Watch"
        startDate="2026-06-23 00:00:00 -0300"
        endDate="2026-06-23 07:00:00 -0300"
        value="HKCategoryValueSleepAnalysisAsleepCore"
    />
</HealthData>
\end{lstlisting}

\section{Testes de funcoes puras}

\begin{lstlisting}[style=devbutterpython]
from src.parser.xml_parser import clean_tag, safe_float


def test_clean_tag_without_namespace():
    assert clean_tag("Record") == "Record"


def test_clean_tag_with_namespace():
    tag = "{http://www.topografix.com/GPX/1/1}trkpt"
    assert clean_tag(tag) == "trkpt"


def test_safe_float_number():
    assert safe_float("90") == 90.0


def test_safe_float_text():
    assert safe_float("HKCategoryValueSleepAnalysisAsleepCore") is None
\end{lstlisting}

\section{Exercicios progressivos}

\subsection{Exercicio 1: Counter}

Conte palavras:

\begin{lstlisting}[style=devbutterpython]
words = ["heart", "steps", "heart", "sleep"]
\end{lstlisting}

Resultado esperado:

\begin{lstlisting}[style=devbutterpython]
{"heart": 2, "steps": 1, "sleep": 1}
\end{lstlisting}

Depois aplique a mesma ideia para contar tipos de \code{Record}.

\subsection{Exercicio 2: batch}

Crie uma funcao que recebe numeros de 1 a 20 e separa em batches de 5.

Entrada:

\begin{lstlisting}[style=devbutterpython]
list(range(1, 21))
\end{lstlisting}

Saida:

\begin{lstlisting}[style=devbutterpython]
[[1, 2, 3, 4, 5], [6, 7, 8, 9, 10], ...]
\end{lstlisting}

\subsection{Exercicio 3: queue}

Crie um producer que coloca numeros em uma fila e um consumer que imprime os numeros.

\subsection{Exercicio 4: threading}

Rode producer e consumer em threads separadas e coloque o nome da thread no log.

\subsection{Exercicio 5: profiling}

Meça o tempo de uma funcao usando \code{perf_counter}.

\subsection{Exercicio 6: banco}

Insira 1000 registros fake no Postgres com:

\begin{itemize}
  \item insert linha por linha;
  \item batch insert;
\end{itemize}

Compare o tempo.

\section{Testes de banco}

No comeco, pode testar manualmente. Depois, use um banco de teste via Docker.

Teste importante:

\begin{itemize}
  \item inserir batch;
  \item contar linhas;
  \item consultar por tipo;
  \item garantir que \code{value_numeric} vira \code{NULL} quando valor e texto.
\end{itemize}

\chapter{Roadmap passo a passo}

\section{Milestone 1: Infra e Postgres}

Entregaveis:

\begin{itemize}
  \item \code{docker-compose.yml};
  \item \code{.env.example};
  \item \code{settings.py};
  \item conexao com Postgres;
  \item schema inicial.
\end{itemize}

Conceitos:

\begin{itemize}
  \item banco relacional;
  \item connection string;
  \item variaveis de ambiente;
  \item tipos SQL;
  \item indices.
\end{itemize}

\section{Milestone 2: Profiler XML}

Objetivo: descobrir o que existe no export.

Entregaveis:

\begin{itemize}
  \item contar tags;
  \item contar tipos de \code{Record};
  \item coletar amostras pequenas;
  \item salvar resultado em JSON ou tabela.
\end{itemize}

Conceitos:

\begin{itemize}
  \item streaming;
  \item \code{Counter};
  \item \code{dict};
  \item \code{set};
  \item O(n).
\end{itemize}

\section{Milestone 3: Parser sincrono}

Entregaveis:

\begin{itemize}
  \item \code{HealthRecord} dataclass;
  \item parser com \code{iterparse};
  \item batch insert;
  \item logs de progresso;
  \item testes com XML pequeno.
\end{itemize}

Conceitos:

\begin{itemize}
  \item dataclass;
  \item list como batch;
  \item memory vs speed;
  \item I/O-bound;
  \item logging.
\end{itemize}

\section{Milestone 4: Benchmark}

Entregaveis:

\begin{itemize}
  \item medir tempo total;
  \item medir records/s;
  \item comparar batch sizes;
  \item salvar resultados em tabela ou markdown.
\end{itemize}

Conceitos:

\begin{itemize}
  \item profiling;
  \item throughput;
  \item fator constante;
  \item tradeoffs.
\end{itemize}

\section{Milestone 5: Queue + Threading}

Entregaveis:

\begin{itemize}
  \item producer;
  \item consumer;
  \item \code{queue.Queue};
  \item backpressure com \code{maxsize};
  \item comparar com pipeline sincrona.
\end{itemize}

Conceitos:

\begin{itemize}
  \item threads;
  \item GIL;
  \item I/O-bound;
  \item concorrencia;
  \item fila FIFO.
\end{itemize}

\section{Milestone 6: GPX concorrente}

Entregaveis:

\begin{itemize}
  \item parser de GPX;
  \item \code{ThreadPoolExecutor};
  \item insert batch de pontos;
  \item benchmark por numero de workers.
\end{itemize}

Conceitos:

\begin{itemize}
  \item fan-out/fan-in;
  \item I/O concorrente;
  \item threads para muitos arquivos.
\end{itemize}

\section{Milestone 7: Multiprocessing}

Entregaveis:

\begin{itemize}
  \item transformacao CPU-bound artificial;
  \item benchmark sync vs threads vs processos;
  \item conclusao sobre GIL.
\end{itemize}

\section{Milestone 8: Insights}

Entregaveis:

\begin{itemize}
  \item batimento medio por dia;
  \item passos por dia;
  \item treinos por semana;
  \item sono por noite;
  \item media movel de 7 dias;
  \item tabela \code{daily_health_summary}.
\end{itemize}

\section{Milestone 9: Dashboard}

Depois que os dados estiverem no banco e os insights basicos existirem, crie uma interface com Streamlit ou FastAPI + frontend.

Ordem recomendada:

\begin{enumerate}
  \item Streamlit para visualizar rapido;
  \item FastAPI se quiser API real;
  \item async depois que houver endpoints I/O-bound claros.
\end{enumerate}

\chapter{Codigo base sugerido para comecar}

Este capitulo junta um esqueleto minimo para iniciar sem perder sua estrutura.

\section{src/app.py}

\begin{lstlisting}[style=devbutterpython]
from pathlib import Path
import psycopg

from src.config.logger import configure_logging
from src.config.settings import Settings
from src.db.repository import HealthRecordRepository
from src.pipelines.sync_pipeline import run_sync_pipeline


def main() -> None:
    configure_logging()
    settings = Settings()

    xml_file = Path("data/export.xml").resolve()

    with psycopg.connect(settings.database_url) as connection:
        repository = HealthRecordRepository(connection)

        run_sync_pipeline(
            xml_file=xml_file,
            repository=repository,
            batch_size=settings.batch_size,
        )


if __name__ == "__main__":
    main()
\end{lstlisting}

\section{Comando para rodar}

\begin{lstlisting}[style=devbutterbash]
poetry run python src/app.py
\end{lstlisting}

\section{Dependencias iniciais}

\begin{lstlisting}[style=devbutterbash]
poetry add "psycopg[binary]"
poetry add --group dev pytest
\end{lstlisting}

Mais tarde:

\begin{lstlisting}[style=devbutterbash]
poetry add pandas streamlit plotly
\end{lstlisting}

\section{Queries iniciais de insight}

\subsection{Tipos mais comuns}

\begin{lstlisting}[style=devbuttersql]
SELECT
    type,
    COUNT(*) AS total
FROM health_records
GROUP BY type
ORDER BY total DESC;
\end{lstlisting}

\subsection{Batimento medio por dia}

\begin{lstlisting}[style=devbuttersql]
SELECT
    DATE(start_date) AS day,
    AVG(value_numeric) AS avg_heart_rate,
    MIN(value_numeric) AS min_heart_rate,
    MAX(value_numeric) AS max_heart_rate,
    COUNT(*) AS samples
FROM health_records
WHERE type = 'HKQuantityTypeIdentifierHeartRate'
GROUP BY DATE(start_date)
ORDER BY day;
\end{lstlisting}

\subsection{Passos por dia}

\begin{lstlisting}[style=devbuttersql]
SELECT
    DATE(start_date) AS day,
    SUM(value_numeric) AS steps
FROM health_records
WHERE type = 'HKQuantityTypeIdentifierStepCount'
GROUP BY DATE(start_date)
ORDER BY day;
\end{lstlisting}

\subsection{Media movel de 7 dias}

\begin{lstlisting}[style=devbuttersql]
WITH daily_steps AS (
    SELECT
        DATE(start_date) AS day,
        SUM(value_numeric) AS steps
    FROM health_records
    WHERE type = 'HKQuantityTypeIdentifierStepCount'
    GROUP BY DATE(start_date)
)
SELECT
    day,
    steps,
    AVG(steps) OVER (
        ORDER BY day
        ROWS BETWEEN 6 PRECEDING AND CURRENT ROW
    ) AS steps_7d_avg
FROM daily_steps
ORDER BY day;
\end{lstlisting}

\section{Comando de validacao com EXPLAIN}

\begin{lstlisting}[style=devbuttersql]
EXPLAIN ANALYZE
SELECT
    DATE(start_date) AS day,
    AVG(value_numeric) AS avg_heart_rate
FROM health_records
WHERE type = 'HKQuantityTypeIdentifierHeartRate'
GROUP BY DATE(start_date)
ORDER BY day;
\end{lstlisting}

Rode antes e depois dos indices para observar o plano de execucao.

\section{Proxima tarefa pratica}

A proxima tarefa deve ser implementar a Milestone 1:

\begin{enumerate}
  \item criar \code{docker-compose.yml};
  \item criar tabelas;
  \item criar \code{Settings};
  \item criar \code{configure_logging};
  \item criar \code{HealthRecord};
  \item criar testes de \code{clean_tag} e \code{safe_float};
  \item rodar um parse pequeno com fixture.
\end{enumerate}

So depois disso processe o arquivo gigante.


\backmatter
\chapter*{Conclusao}
Este documento transforma seu parser do Apple Health em um laboratorio pratico de engenharia: primeiro funciona, depois mede, depois otimiza e compara. A implementacao final deve permitir gerar insights pessoais a partir de dados reais, enquanto voce treina streaming, filas, Postgres, logging, profiling, concorrencia e estruturas de dados em Python.

\end{document}
